\documentclass[11pt,a4paper]{article}
\usepackage[margin=2.2cm]{geometry}
\usepackage[T1]{fontenc}
\usepackage[utf8]{inputenc}
\usepackage{xcolor,listings,hyperref,enumitem}
\definecolor{navy}{HTML}{17365D}
\definecolor{pale}{HTML}{F3F6FA}
\lstdefinestyle{code}{basicstyle=\ttfamily\small,breaklines=true,frame=single,backgroundcolor=\color{pale},rulecolor=\color{gray!35},numbers=left,numberstyle=\tiny\color{gray},numbersep=8pt,showstringspaces=false,tabsize=4,columns=fullflexible,keepspaces=true}
\lstset{style=code}
\hypersetup{colorlinks=true,urlcolor=navy}
\title{\textcolor{navy}{Task 1: Review of Programming Basics}\\\large Python and C\# Console Applications}
\author{Programming Exercises -- Complete Solutions}
\date{November 2026}
\begin{document}\maketitle
\section*{Overview}
This report presents working console-based solutions for three exercises in both Python and C\#. Each solution reads input, computes a result, and prints it to the console. Input validation prevents invalid operations and grades.
\section{Simple Two-Number Calculator}
\textbf{Goal:} Read two numbers and an operator (+, -, *, /), then display the result. Division by zero is handled explicitly.
\subsection*{Python}
\begin{lstlisting}[language=Python]
try:
    first = float(input("First number: "))
    second = float(input("Second number: "))
    operator = input("Operation (+, -, *, /): ").strip()

    if operator == "+":
        result = first + second
    elif operator == "-":
        result = first - second
    elif operator == "*":
        result = first * second
    elif operator == "/":
        if second == 0:
            raise ZeroDivisionError("Cannot divide by zero.")
        result = first / second
    else:
        raise ValueError("Unknown operation.")

    print(f"Result: {result:g}")
except (ValueError, ZeroDivisionError) as error:
    print(f"Error: {error}")
\end{lstlisting}
\subsection*{C\#}
\begin{lstlisting}
using System;
using System.Globalization;

Console.Write("First number: ");
bool firstOk = double.TryParse(Console.ReadLine(),
    NumberStyles.Float, CultureInfo.InvariantCulture, out double a);
Console.Write("Second number: ");
bool secondOk = double.TryParse(Console.ReadLine(),
    NumberStyles.Float, CultureInfo.InvariantCulture, out double b);
Console.Write("Operation (+, -, *, /): ");
string op = Console.ReadLine()?.Trim() ?? "";

if (!firstOk || !secondOk)
{
    Console.WriteLine("Error: Invalid number.");
}
else if (op == "/" && b == 0)
{
    Console.WriteLine("Error: Cannot divide by zero.");
}
else
{
    double result;
    if (op == "+") result = a + b;
    else if (op == "-") result = a - b;
    else if (op == "*") result = a * b;
    else if (op == "/") result = a / b;
    else
    {
        Console.WriteLine("Error: Unknown operation.");
        return;
    }
    Console.WriteLine($"Result: {result:G}");
}
\end{lstlisting}
\textbf{Example:} Inputs: \texttt{10}, \texttt{20}, \texttt{+}. Output: \texttt{Result: 30}.
\section{Temperature Converter}
\textbf{Goal:} Select C to convert Celsius to Fahrenheit, or F to convert Fahrenheit to Celsius. The formulas are
\[F = 1.8C + 32, \qquad C = \frac{F-32}{1.8}.\]
\subsection*{Python}
\begin{lstlisting}[language=Python]
try:
    direction = input("Convert from C or F? ").strip().upper()
    if direction not in ("C", "F"):
        raise ValueError("Choose C or F.")
    temperature = float(input("Temperature: "))

    if direction == "C":
        converted = temperature * 1.8 + 32
        print(f"{temperature:g} C = {converted:g} F")
    else:
        converted = (temperature - 32) / 1.8
        print(f"{temperature:g} F = {converted:g} C")
except ValueError as error:
    print(f"Error: {error}")
\end{lstlisting}
\subsection*{C\#}
\begin{lstlisting}
using System;
using System.Globalization;

Console.Write("Convert from C or F? ");
string direction = (Console.ReadLine() ?? "").Trim().ToUpper();
if (direction != "C" && direction != "F")
{
    Console.WriteLine("Error: Choose C or F.");
    return;
}
Console.Write("Temperature: ");
if (!double.TryParse(Console.ReadLine(), NumberStyles.Float,
    CultureInfo.InvariantCulture, out double temperature))
{
    Console.WriteLine("Error: Invalid temperature.");
    return;
}
if (direction == "C")
{
    double fahrenheit = temperature * 1.8 + 32;
    Console.WriteLine($"{temperature:G} C = {fahrenheit:G} F");
}
else
{
    double celsius = (temperature - 32) / 1.8;
    Console.WriteLine($"{temperature:G} F = {celsius:G} C");
}
\end{lstlisting}
\textbf{Example:} Inputs: \texttt{C}, \texttt{0}. Output: \texttt{0 C = 32 F}.
\section{Student Grade Point Average}
\textbf{Goal:} Read a positive number of grades on the 1--6 scale, calculate their arithmetic mean, and print whether the student passed (average at least 3.0).
\subsection*{Python}
\begin{lstlisting}[language=Python]
try:
    count = int(input("Number of grades: "))
    if count <= 0:
        raise ValueError("Enter at least one grade.")

    total = 0.0
    for index in range(count):
        grade = float(input(f"Grade {index + 1} (1-6): "))
        if not 1 <= grade <= 6:
            raise ValueError("Grades must be between 1 and 6.")
        total += grade

    average = total / count
    print(f"Average: {average:.2f}")
    if average >= 3.0:
        print("The student passed.")
    else:
        print("The student did not pass.")
except ValueError as error:
    print(f"Error: {error}")
\end{lstlisting}
\subsection*{C\#}
\begin{lstlisting}
using System;
using System.Globalization;

Console.Write("Number of grades: ");
if (!int.TryParse(Console.ReadLine(), out int count) || count <= 0)
{
    Console.WriteLine("Error: Enter a positive number of grades.");
    return;
}
double total = 0;
for (int i = 0; i < count; i++)
{
    Console.Write($"Grade {i + 1} (1-6): ");
    if (!double.TryParse(Console.ReadLine(), NumberStyles.Float,
        CultureInfo.InvariantCulture, out double grade) ||
        grade < 1 || grade > 6)
    {
        Console.WriteLine("Error: Grades must be between 1 and 6.");
        return;
    }
    total += grade;
}
double average = total / count;
Console.WriteLine($"Average: {average:F2}");
if (average >= 3.0)
    Console.WriteLine("The student passed.");
else
    Console.WriteLine("The student did not pass.");
\end{lstlisting}
\textbf{Example:} Inputs: \texttt{4}, \texttt{5}, \texttt{3}, \texttt{4}, \texttt{2}. Output: \texttt{Average: 3.50} followed by \texttt{The student passed.}
\section*{Execution and Submission}
Each code sample is a standalone console application. Save Python examples as \texttt{task1.py}, \texttt{task2.py}, and \texttt{task3.py}, and run with \texttt{python task1.py} (similarly for the other tasks). For C\#, create a separate console project for each task using \texttt{dotnet new console}, replace its \texttt{Program.cs} with the corresponding code, and execute \texttt{dotnet run}. C\# examples use top-level statements (C\# 9 or later). Numeric input in C\# uses a dot as decimal separator.

For submission, push the source files to GitHub, preferably with separate Python and C\# directories and a README containing run instructions. Use a feature branch (for example, \texttt{feature/programming-basics}) and merge it into \texttt{main} after testing. A GitHub repository link must be created and submitted separately; this document does not represent a published repository.
\end{document}
